% =============================================================
% Capitulo 5 - Control automatico e interfaces humano-maquina
% =============================================================
\chapter{Control autom\'atico e interfaces humano-m\'aquina}
\label{cap:control_hmi}

\section{Lazo de control de sintonizaci\'on de longitud de onda}
\label{sec:lazo_control}

El lazo de control abordado en este proyecto tiene como variable
controlada la temperatura de cavidad del l\'aser seeder, de la cual
depende proporcionalmente la longitud de onda sintonizada, y como
variable medida la diferencia entre las amplitudes de las se\~nales
$P_+$ y $P_-$ obtenidas de los fotodetectores del sistema (ver
Secci\'on~\ref{sec:hsrl_pilar}) \cite{sat2026}. El sistema encuentra el
punto de trabajo correcto cuando dicha diferencia se anula, es decir,
cuando las se\~nales $P_+$ y $P_-$ se igualan \cite{sat2026}.

\section{Control proporcional aplicado al error de sintonizaci\'on}
\label{sec:control_proporcional}

La se\~nal de error del lazo se define como

\begin{equation}
  e(k) = P_+(k) - P_-(k)
  \label{eq:error_control}
\end{equation}

\noindent donde $P_+(k)$ y $P_-(k)$ son los valores de pico medidos por
los detectores en la iteraci\'on $k$ \cite{informeMarzo2026}. Se
utiliza una ley de control proporcional para actualizar la temperatura
de consigna del seeder:

\begin{equation}
  T(k) = T(k-1) - K_p \, e(k)
  \label{eq:ley_control}
\end{equation}

\noindent siendo $K_p$ la constante proporcional del controlador
\cite{informeMarzo2026}. El detalle de la implementaci\'on de esta ley
de control en el firmware del microcontrolador se presenta en la
Secci\'on~\ref{sec:algoritmo_control}.

% PLACEHOLDER: justificar teoricamente la eleccion de un controlador
% proporcional (sin terminos integral/derivativo) y, de corresponder,
% incorporar el analisis de estabilidad y el criterio de ajuste de
% Kp una vez determinado experimentalmente.

\section{Principios de dise\~no de interfaces humano-m\'aquina (norma ISA-101)}
\label{sec:isa101}

El software de interfaz humano-m\'aquina (HMI) del sistema se dise\~na
tomando como base los lineamientos de la norma ISA-101
\cite{informeAbril2026}. Dicha interfaz, desarrollada en Python, tiene
por objetivo permitir la interacci\'on entre el operador y el
microcontrolador de la placa de control, incluyendo la visualizaci\'on
del estado de sintonizaci\'on del sistema y la modificaci\'on de
par\'ametros y funciones de operaci\'on \cite{sat2026,informeAbril2026}.

% PLACEHOLDER: desarrollar los principios generales de la norma ISA-101
% (jerarquia de pantallas, uso del color, gestion de alarmas) y su
% aplicacion concreta al diseno del HMI una vez avanzado su desarrollo
% (ver Capitulo \ref{cap:software_hmi}).
